\chapter{Một từ cho tâm trạng}
\begin{tinhhuongbox}{Tình huống | Situation}
\songngu{Sau kiểm tra, lớp có thể vừa mệt vừa tiếc vừa buồn cười vì một câu mình viết sai.}{After a test, the class can feel tired, regretful, and oddly amused by one answer they got wrong.}
\songngu{Nếu không xả một chút, tiết sau sẽ phải gánh hết phần dư âm đó.}{If you do not release some of it, the next lesson will carry all that leftover emotion.}
\end{tinhhuongbox}
\begin{noitrenlop}
\songngu{Mỗi em một từ cho tâm trạng sau kiểm tra thôi, đừng kể nguyên series dài tập.}{One word each for your post-test mood, please. Do not tell the whole long-running series.}
\end{noitrenlop}
\begin{vihieuqua}
\songngu{Một từ thì xả được mà chưa làm loãng tiết học kế tiếp.}{One word lets students release tension without watering down the next lesson.}
\end{vihieuqua}
\begin{englishpocket}
\songngu{Cụm nên nhớ: \textit{Post-test mood.}}{Phrase to keep: \textit{Post-test mood.}}
\end{englishpocket}
\chapter{Phần cũ dừng ở đây}
\begin{tinhhuongbox}{Tình huống | Situation}
\songngu{Nhiều lớp không tự chuyển trạng thái được sau kiểm tra.}{Many classes cannot switch states by themselves after a test.}
\songngu{Người dạy phải vạch một đường ranh rõ cho các em bước qua.}{The teacher has to draw a clear line for students to step across.}
\end{tinhhuongbox}
\begin{noitrenlop}
\songngu{Rồi, phần lo lắng dừng ở đây, phần học tiếp bắt đầu từ đây.}{All right, the worrying stops here, and the next part of learning starts here.}
\end{noitrenlop}
\begin{vihieuqua}
\songngu{Một câu tách nhịp rõ ràng làm lớp nhẹ đi hẳn.}{A clear transition sentence makes the class feel noticeably lighter.}
\end{vihieuqua}
\begin{englishpocket}
\songngu{Cụm nên nhớ: \textit{The worrying stops here.}}{Phrase to keep: \textit{The worrying stops here.}}
\end{englishpocket}
\chapter{Cho lớp thắng một câu dễ}
\begin{tinhhuongbox}{Tình huống | Situation}
\songngu{Sau kiểm tra, học sinh rất cần cảm giác mình vẫn làm được điều gì đó.}{After a test, students badly need the feeling that they can still do something well.}
\songngu{Một chiến thắng nhỏ đầu tiên có tác dụng hồi năng lượng rất mạnh.}{A small early win can restore energy very strongly.}
\end{tinhhuongbox}
\begin{noitrenlop}
\songngu{Câu đầu tiên tôi cho dễ một chút, để lớp nhớ rằng não mình chưa nghỉ hưu.}{I will make the first question a little easier, so the class remembers that your brains have not retired.}
\end{noitrenlop}
\begin{vihieuqua}
\songngu{Lớp bật cười rồi bước vào phần mới với ít sợ hơn.}{The class laughs and then enters the new part with less fear.}
\end{vihieuqua}
\begin{englishpocket}
\songngu{Cụm nên nhớ: \textit{Your brains have not retired.}}{Phrase to keep: \textit{Your brains have not retired.}}
\end{englishpocket}
\chapter{Đặt tên cho bài kiểm tra vừa rồi}
\begin{tinhhuongbox}{Tình huống | Situation}
\songngu{Đặt tên vui cho một chuyện căng là cách làm nó nhẹ hơn mà không phủ nhận nó.}{Giving a stressful event a funny title makes it lighter without denying it.}
\songngu{Học sinh thường rất thích trò này.}{Students usually love this move.}
\end{tinhhuongbox}
\begin{noitrenlop}
\songngu{Nếu đặt tên cho bài kiểm tra vừa rồi như tên phim, lớp mình đặt là gì?}{If we had to give that test a movie title, what would we call it?}
\end{noitrenlop}
\begin{vihieuqua}
\songngu{Tiếng cười nhỏ lúc này giúp cắt độ căng đang bám trong lớp.}{A small laugh at this moment helps cut through the tension still hanging in the room.}
\end{vihieuqua}
\begin{englishpocket}
\songngu{Cụm nên nhớ: \textit{What would we call it?}}{Phrase to keep: \textit{What would we call it?}}
\end{englishpocket}
\chapter{Đừng mở bài mới bằng mặt cũ}
\begin{tinhhuongbox}{Tình huống | Situation}
\songngu{Có khi nội dung mới không nặng, chỉ là gương mặt người học còn nặng.}{Sometimes the new content is not heavy; the students' faces are.}
\songngu{Đầu vào của bài mới phải nhẹ hơn một nấc.}{The entrance to the new lesson should be one step lighter.}
\end{tinhhuongbox}
\begin{noitrenlop}
\songngu{Giờ mình vào phần mới bằng mặt tươi hơn một chút nhé, đừng để đề kiểm tra ngồi học tiếp với mình.}{Let us enter the new part with slightly fresher faces. Do not let the test come to class with us.}
\end{noitrenlop}
\begin{vihieuqua}
\songngu{Câu này thừa nhận cảm xúc cũ nhưng không cho nó cầm lái.}{This line acknowledges the old feeling without letting it drive.}
\end{vihieuqua}
\begin{englishpocket}
\songngu{Cụm nên nhớ: \textit{Do not let the test come to class with us.}}{Phrase to keep: \textit{Do not let the test come to class with us.}}
\end{englishpocket}